\uuid{l7rf}
\titre{Matrice/Rang}
\auteur{hanani ; mzari}
\organisation{exo7}
\licence{CC BY-NC-SA 4.0}
\datecreate{2020-01-20}
\niveau{L1}
\module{Algèbre}
\chapitre{Matrice}
\sousChapitre{Propriétés élémentaires, généralités}
\difficulte{2}
\theme{Calcul matriciel, Matrice, Rang}
\exo7qcmid{lille-2.matrices.009}

\contenu{
\begin{question}
\qclassification{108.01}
\qtags{motcle=Matrice/Rang}
On considère les matrices :
$$A=  \left(\begin{array}{rcc}
1&0&1\\\end{array}\right),\quad B=
\left(\begin{array}{rcc}2&-4\\
1&-2\\ 0&0\\
\end{array}\right),\quad C=
\left(\begin{array}{rcc}1&0&0\\1&1&1\\ 0&1&2\\
\end{array}\right)\quad \mbox{et}\quad D=
\left(\begin{array}{rcc}1&0&1\\1&1&0\\ 0&1&-1\\
\end{array}\right). $$
Quelles sont les assertions vraies ?
\begin{answers}
\bad{Le rang de $A$ est $3$}.
\good{Le rang de $B$ est $1$}.
\good{Le rang de $C$ est $3$}.
\bad{Le rang de $D$ est $3$}.
\end{answers}
\begin{explanations} Le rang d'une matrice est le nombre maximum de vecteurs colonnes ou lignes qui sont linéairement indépendants. Le rang de $A$ est $1$, le rang de $B$ est $1$, le rang de $C$ est $3$ et le rang de $D$ est $2$.
\end{explanations}
\end{question}
}
